\documentclass[12pt,a4paper]{article}
\usepackage[utf8]{inputenc}
\usepackage{url}
\usepackage{hyperref}
\usepackage{amsmath}
\usepackage{amsfonts}
\usepackage{amssymb}
\usepackage{graphicx}
\usepackage{float}
\usepackage{lipsum}
\usepackage{titlesec}
%\usepackage[font=scriptsize,labelfont=bf]{caption} % for caption font size
%\usepackage{siunitx} %for si units shortcuts
\usepackage[left=1.1cm, right=1.1cm, top=2.00cm, bottom=2.00cm]{geometry}
%\usepackage{subcaption}

%\begin{multicols}{2}
%\end{multicols}
%definecolor{azul}{rgb}{0.0, 0.53, 0.74}

% Formattazione sezioni
%\titleformat{\section}{\large\bfseries}{\thesection.}{0.5em}{}
%\titleformat{\subsection}{\normalsize\bfseries}{\thesubsection.}{0.5em}{}
%\renewcommand{\abstractname}{\large Abstract} % Cambia la dimen 

\begin{document}

\begin{titlepage}
    \centering
    \vspace*{2cm} % Spazio verticale per centrare
    
    {\huge \textbf{ Atoms in DFT Report}}
    \vspace*{1,5cm} % Spazio verticale per centrare
    
    {\large Univerity of Bologna , Physics , Computational Physics Laboratory }
    \vspace*{1cm} % Spazio verticale per centrare

    {\large Bertasi Leonardo, Mulazzani Marco e Hutifa }
    \vspace*{1cm} % Spazio verticale per centrare
    
     \large Date: 30/01/2026

    %\includegraphics[width=0.4\textwidth]{logo.png} % Logo opzionale\\[1.5cm]
   		

    %\begin{abstract}
    %\end{abstract}

    \vfill  % Riempie lo spazio verticale per centrare il contenuto
\end{titlepage}



\section{Introduction}
The first part of the lab work on computational physics is focus on the
Calculation of Electronic Energy levels of differnt atoms such as 
Helium, carbon, Nitrogen and oxygen with VASP.
This type of software takes the calculations inputs from a few standard
 files : 
 \begin{itemize}
    \item POSCAR : contains the information about the lattice 
    and the position of the atoms in the unit cell.
    \item POTCAR : contains the information about the 
    pseudopotential of the atoms in the system.
    \item INCAR : contains the information about the settings 
    of the calculation such as the type
    \item KPOINTS : contains the information about the k-points
     grid used for the calculation.
    \item OUTCAR : contains the output of the calculation such
     as the energy levels and the forces on the atoms.
 \end{itemize}

 Fot this calculation these are the only required even if there 
 are other files that can be used for more specific calculations.


\section{Calculations Method}
After affect test about He atom with a fixed lattice parameter,
 


we performed the calculations for 
the other atoms with the same settings.


\begin{figure}[h]
    \centering
        \includegraphics[scale=0.6 ]{ char.png}
        \caption{Resonance Characteristic Curve of the NCHR cantilever. Notice the unusual 
        difference in the phase plot, an error must have occured during data acquisition. 
        Amplitude modulus set to 20nm, set point 15.1nm.}
        \label{fig:eg}
\end{figure}    


\section{Observations and Results}
 
 

 \begin{table}[h!]
    
    \centering
    \begin{tabular}{|c|c|c|c|c|} % Column specification with vertical lines
        \hline % Horizontal line
         \textbf{Key Parameters} & Image Size  & Scan Rate & Set Point & Drive Amplitude \\
        \hline
          & 15x15 $\mu m$ & 0.5 Hz & 15.1 nm  & 20 nm \\
        \hline
    \end{tabular}       
    \caption{Main parameters used during NCM-AFM scan are shown here. Some info could 
     be missing due to data error explained before.}
    \label{tab:afm_modes}
\end{table}



\section{Final Results}
While the experiment seemed to be performed without any important issues during the AFM operation and data acquisition in the laboratory, the afterwards analysis of the collected data was quite problematic.
On a good note the available results were as predicted, however the details were 
insufficient to pursue a complete evaluation. \\
Results for the resonance curve were consistent for the amplitude but a significant shift was noticed on the phase curve. This can be due to different reasons, with the most probable one being bad settings of the software during acquisition.\\
The results of the NCM-AFM imaging were fine but no height scale bar was acquired. Finally, for the AFM Spectroscopy the results were as predicted from
the theoretical curves with the Attractive regime approximation.  



\end{document}


\begin{figure}[h]
    \centering
        \begin{subfigure}{0.48 \textwidth}
            \centering
           \includegraphics[width= 0.9 \linewidth]{afm_setup_lockin.png}
           \vspace{0.2cm} 
           \caption{ A schematic view of the AFM experimental setup, including the lock-in amplifier and function generator. 
            An additional function generator is used to control the sample 2D's scan \cite{cramer_notes}.}
            \label{fig:setting}
        \end{subfigure}
        \hfill
        \begin{subfigure}{0.48 \textwidth}      
            \centering
           \includegraphics[width =  0.8 \linewidth]{afm_modes.png}
            \caption{  } 
             \label{fig:setup}
        \end{subfigure}
  
\end{figure}    
 